(1) $L=m\ell^2(\dot\theta^2+\sin^2\theta\dot\phi^2)/2+mg\ell\cos\theta$；方位角是循环坐标，故 $L_z=m\ell^2\sin^2\theta\dot\phi$ 守恒。该守恒量不能误当成绕绳方向的角动量。

(2) 消去 $\dot\phi=L_z/(m\ell^2\sin^2\theta)$ 后 $E=m\ell^2\dot\theta^2/2+U_{\rm eff}$，$U_{\rm eff}=L_z^2/(2m\ell^2\sin^2\theta)-mg\ell\cos\theta$。在 $L_z\ne0$ 时极点有无限势垒，允许的角满足 $E\ge U_{\rm eff}(\theta)$。

(3) 驻点要求 $L_z^2\cos\theta_0=m^2g\ell^3\sin^4\theta_0$。使用竖直角动量表达式得到 $\dot\phi_0^2=g/(\ell\cos\theta_0)$。由于左式非负，圆锥平衡要求 $0<\theta_0<\pi/2$，下半球的延拓并非物理解。

(4) 求二阶导并代入平衡等式可得 $\omega_{\rm nut}^2=\dot\phi_0^2(1+3\cos^2\theta_0)$。该频率比介于1与2之间；通常不与进动频率整比，因此球面摆在实验室系的轨道一般呈章动而非严格闭合圆。
